\chapter{Implementation}
\label{chap:implementation}

\autoref{chap:sysmodel} defined the crash-aware system at a conceptual level: a decentralized IEEE~802.11p highway network, a synthetic two-class \gls{be}/\gls{vo} traffic model, five \gls{mac} policies with a local three-mode controller, and a three-phase crash timeline.
This chapter maps that model onto the concrete software artifact developed for this dissertation, so that a reader who already understands the design can reproduce the build, follow the packet path, and locate every study-specific module without consulting undocumented source comments.

The realization is an \gls{omnet}~6.1~\cite{varga2008omnetpp} / \gls{inet}~4.5.4~\cite{inet_framework} / Veins~5.3.1~\cite{veins_sommer2011} project whose source, scenario packages, and analysis tooling are publicly available~\cite{aguiar2026_veins_qos}.
Road mobility is provided by \gls{sumo} coupled through \gls{traci}~\cite{lopez2018sumo}.
The chapter is organized as follows.
Section~\ref{sec:impl:toolchain} pins the toolchain and repository layout.
Section~\ref{sec:impl:host} describes the vehicle host and scenario compound modules.
Section~\ref{sec:impl:path} traces the end-to-end packet path.
Section~\ref{sec:impl:apps} details the application modules.
Section~\ref{sec:impl:classifier} covers the DSCP-to-Voice classifier.
Section~\ref{sec:impl:mac} explains the instrumented \gls{mac} and the adaptive \gls{hcf} controller.
Section~\ref{sec:impl:policies} documents how the five policies are wired in configuration.
Section~\ref{sec:impl:scenarios} presents the highway scenario packages, the consolidated simulation-parameter table, and the configuration matrix.
Section~\ref{sec:impl:instrumentation} lists recorded signals and the analysis tooling surface.
Section~\ref{sec:impl:scope} states implementation scope and deliberate omissions.
Seeds, statistical aggregation, and measured results belong to \autoref{chap:evaluation} and are not reported here.

%%%%%%%%%%%%%%%%%%%%%%%%%%%%%%%%%%%%%%%%%%%%%%%%%%%%%%%%%%%%%%%%%%%%%%%%%%%%%%%%
\section{Simulation Toolchain and Repository Layout}
\label{sec:impl:toolchain}

The Background chapter motivates bidirectional road-and-network coupling for inter-vehicle protocol studies (\autoref{sec:bg:sim}).
Here the same stack is pinned to the versions used for the experiments and narrowed to the modules that realize the System Model.

\subsection{Framework versions}
\label{sec:impl:versions}

\gls{omnet}~6.1 provides the discrete-event engine, component model, and \texttt{.ini}/\texttt{.ned} configuration system in which every vehicle node is assembled~\cite{varga2008omnetpp}.
\gls{inet}~4.5.4 supplies the Internet protocol stack, IEEE~802.11p-capable \gls{mac} and \gls{phy}, radio propagation abstractions, and the Hybrid Coordination Function (\gls{hcf}) into which the study's adaptive controller is substituted for selected policies~\cite{inet_framework,IEEE80211e_2005}.
Veins couples \gls{omnet}/\gls{inet} to \gls{sumo} through \gls{traci}, so highway mobility, lane discipline, and the crash-node stop command remain consistent with the road microsimulation~\cite{veins_sommer2011,lopez2018sumo}.
In the artifact tree, \texttt{inet/} and \texttt{veins/} are git submodules consumed as read-only frameworks, pinned to the \gls{inet}~4.5.4 and Veins~5.3.1 releases; the study-authored code lives under \texttt{veins\_qos/}~\cite{aguiar2026_veins_qos}.

\subsection{Repository layout}
\label{sec:impl:layout}

\autoref{tab:impl:layout} summarizes the directories that matter for implementation and reproducibility.
Only the active highway packages participate in the dissertation experiment matrix; older square and early highway packages remain in the tree as development history and are not part of the reported study~\cite{aguiar2026_veins_qos}.

\begin{table}[htb]
	\ABNTEXfontereduzida
	\caption{\label{tab:impl:layout}Artifact directory layout relevant to the crash-aware implementation.}
	\begin{center}
		\begin{tabularx}{\linewidth}{@{}>{\ttfamily\raggedright\arraybackslash}p{4.4cm} X@{}}
			\toprule
			\textbf{Path} & \textbf{Role} \\
			\midrule
			veins\_qos/src/traffic/ &
			  Study applications: \texttt{CritPacketSender}, \texttt{CrashBurstApp}. \\
			\addlinespace
			veins\_qos/src/qos/ &
			  Explicit DSCP-to-user-priority classifier. \\
			\addlinespace
			veins\_qos/src/mac/ &
			  Instrumented IEEE~802.11 \gls{mac}, adaptive \gls{hcf}, and \gls{fsm} controller. \\
			\addlinespace
			veins\_qos/src/veins\_inet/ &
			  Veins$\leftrightarrow$\gls{inet} host, mobility, manager, and UDP application base. \\
			\addlinespace
			veins\_qos/simulations/veins\_inet\_highway\_\{light,heavy\}/ &
			  Active density packages: SUMO network, obstacles, \texttt{omnetpp.ini}, runners. \\
			\addlinespace
			kpi\_dashboard/ &
			  Offline parser of OMNeT++ \texttt{.sca}/\texttt{.vec} result files and figure export. \\
			\addlinespace
			inet/, veins/ &
			  Framework submodules (read-only). \\
			\bottomrule
		\end{tabularx}
	\end{center}
	\fonte{Prepared by the author.}
\end{table}

\subsection{Attribution}
\label{sec:impl:attribution}

The Veins/\gls{inet} integration modules under \texttt{veins\_qos/src/veins\_inet/} derive from the Veins project (Sommer et al.)~\cite{veins_sommer2011}.
The \texttt{VeinsInetTransparentMobility} component incorporates work by Jannusch Bigge (TU Dresden); it is present in the tree but is \emph{not} selected by the active highway configurations, which use \texttt{VeinsInetMobility} instead~\cite{aguiar2026_veins_qos}.
Study-specific logic is confined to the application generators, the classifier, the instrumented \gls{mac}, and---for the adaptive policies only---a replacement \gls{hcf} submodule, so that physical-layer connectivity, mobility, and offered-load profiles stay shared across all five \gls{mac} policies (\autoref{sec:sys:network}).

%%%%%%%%%%%%%%%%%%%%%%%%%%%%%%%%%%%%%%%%%%%%%%%%%%%%%%%%%%%%%%%%%%%%%%%%%%%%%%%%
\section{Host and Scenario Compound Modules}
\label{sec:impl:host}

\subsection{Vehicle host}
\label{sec:impl:car}

Each simulated vehicle is an instance of \texttt{veins\_qos.veins\_inet.VeinsInetCar}, a thin extension of INET's \texttt{AdhocHost} that carries a single IEEE~802.11 wireless interface (\texttt{wlan0}), an IPv4 stack with \texttt{HostAutoConfigurator}, and a Veins mobility submodule~\cite{inet_framework,aguiar2026_veins_qos}.
The manager launches cars with

\begin{quote}
\ttfamily *.manager.moduleType = "veins\_qos.veins\_inet.VeinsInetCar"
\end{quote}

so that every TraCI-injected node receives the same compound type.
Mobility is bound to \texttt{veins\_qos.veins\_inet.VeinsInetMobility}, which mirrors SUMO positions, headings, and presence into the INET mobility model~\cite{veins_sommer2011}.
IPv4 autoconfiguration joins the multicast group \texttt{224.0.0.1} on \texttt{wlan0}, matching the application destination used by both traffic generators.

\subsection{Application base}
\label{sec:impl:appbase}

Both study applications extend \texttt{VeinsInetApplicationBase}, itself derived from INET's \texttt{ApplicationBase} with Veins-aware \gls{udp} socket setup~\cite{inet_framework,aguiar2026_veins_qos}.
The base class opens a \gls{udp} socket on \texttt{wlan0}, disables multicast loopback (\texttt{setMulticastLoop(false)}), and provides helpers for timed transmission and TraCI access.
Study applications therefore inherit a common multicast send path and concentrate their own logic on marking, timing, and instrumentation.

\subsection{Scenario network}
\label{sec:impl:scenario-ned}

Each active package defines a \texttt{Scenario} network \gls{ned} that composes:

\begin{itemize}
  \item \texttt{Ieee80211DimensionalRadioMedium} --- shared radio medium for all nodes;
  \item \texttt{VeinsInetManager} --- TraCI scenario manager that launches SUMO through \texttt{veins\_launchd} on \texttt{localhost:9999} using \texttt{highway.launchd.xml};
  \item \texttt{PhysicalEnvironment} --- obstacle geometry from \texttt{highway.obstacles.xml};
  \item visualizers (optional for interactive runs);
  \item an initial \texttt{VeinsInetCar} placeholder; additional cars are injected by TraCI as SUMO vehicles appear.
\end{itemize}

This compound is the OMNeT++ \texttt{network} selected by every runnable configuration in the package~\cite{aguiar2026_veins_qos}.

%%%%%%%%%%%%%%%%%%%%%%%%%%%%%%%%%%%%%%%%%%%%%%%%%%%%%%%%%%%%%%%%%%%%%%%%%%%%%%%%
\section{End-to-End Packet Path}
\label{sec:impl:path}

\autoref{fig:impl:stack} and the steps below summarize how a packet moves from creation to KPI recording.
The path is identical for Best Effort and Voice except for DSCP marking, classifier mapping, and---under adaptive policies---\gls{hcf}-level Best Effort gating.

\begin{figure}[htb]
	\caption{\label{fig:impl:stack}Implemented cross-layer path from study applications to the shared radio. Adaptive \gls{hcf} modules appear only under the \texttt{edca\_v2x\_vo\_*} policies.}
	\begin{center}
	\ABNTEXfontereduzida
	\begin{tabularx}{\linewidth}{>{\centering\arraybackslash}X}
		\toprule
		\textbf{Application:} \texttt{CritPacketSender} (\gls{dscp}~0) or \texttt{CrashBurstApp} (\gls{dscp}~46); \texttt{CreationTimeTag}, optional sequence tag \\
		\midrule
		$\downarrow$ \\
		\textbf{UDP / IPv4:} multicast to \texttt{224.0.0.1} on \texttt{wlan0} (Crash destPort~9001; Crit local/dest~9001) \\
		\midrule
		$\downarrow$ \\
		\textbf{Classifier (EDCA policies):} \texttt{QosClassifier} maps \gls{dscp}~46~$\rightarrow$~\gls{up} Voice; others~$\rightarrow$~\gls{up}~0 \\
		\midrule
		$\downarrow$ \\
		\textbf{\gls{mac}:} \texttt{V2xIeee80211Mac}; \texttt{plain} uses \gls{dcf}; EDCA uses stock \texttt{Hcf} or \texttt{V2xHcf} \\
		\midrule
		$\downarrow$ \\
		\textbf{\gls{phy}:} \texttt{Ieee80211DimensionalRadio}, free-space path loss, \texttt{IdealObstacleLoss}, static \gls{snir} \\
		\midrule
		$\downarrow$ \\
		\textbf{Reception KPIs:} \texttt{CritPacketSender} on every node (port~9001) records BE/VO RX and delay \\
		\bottomrule
	\end{tabularx}
	\end{center}
	\fonte{Prepared by the author.}
\end{figure}

\begin{enumerate}
  \item The sending application creates a UDP payload, attaches a \texttt{CreationTimeTag} (and, for crash bursts, a sequence request), and marks the packet with a \texttt{DscpReq} of~0 or~46.
  \item \texttt{VeinsInetApplicationBase::sendPacket} delivers the frame to UDP, which sends to multicast \texttt{224.0.0.1}.
  \item IPv4 forwards onto \texttt{wlan0}. Under EDCA policies, \texttt{QosClassifier} converts DSCP into a \texttt{UserPriorityReq} before the frame enters the IEEE~802.11 \gls{mac}~\cite{RFC8325_WiFi_QoS}.
  \item \texttt{V2xIeee80211Mac} either enqueues into a single \gls{dcf} pending queue (\texttt{plain}) or into \gls{edca} access-category queues (\texttt{edca\_*}). Adaptive policies consult \texttt{V2xEdcaFsmController} before granting Best Effort channel access~\cite{IEEE80211e_2005}.
  \item The dimensional 802.11p radio transmits on the shared medium; receivers reverse the path through \gls{phy}, \gls{mac}, IP, and UDP.
  \item Every node's \texttt{CritPacketSender} (\texttt{app[0]}, port~9001) processes received datagrams, classifies them by DSCP into BE or VO KPI streams, skips self-originated addresses, and optionally deduplicates Voice repeats.
\end{enumerate}

Crash-alert frames are therefore generated by \texttt{CrashBurstApp} but counted at receivers by \texttt{CritPacketSender}, because the crash application sends to destination port~9001---the same port on which every vehicle listens for cooperative traffic~\cite{aguiar2026_veins_qos}.

%%%%%%%%%%%%%%%%%%%%%%%%%%%%%%%%%%%%%%%%%%%%%%%%%%%%%%%%%%%%%%%%%%%%%%%%%%%%%%%%
\section{Application Layer}
\label{sec:impl:apps}

Two custom application modules implement the synthetic traffic model of \autoref{sec:sys:traffic}.
Both are instantiated on every vehicle (\texttt{numApps = 2}); only the configured crash target activates the Voice overlay.

\subsection{\texttt{CritPacketSender}: periodic Best Effort}
\label{sec:impl:crit}

\texttt{veins\_qos.traffic.CritPacketSender} extends \texttt{VeinsInetApplicationBase} and generates the ordinary background flow~\cite{aguiar2026_veins_qos}.
\autoref{tab:impl:crit} lists the principal NED parameters and recorded statistics.

\begin{table}[htb]
	\ABNTEXfontereduzida
	\caption{\label{tab:impl:crit}Principal parameters and statistics of \texttt{CritPacketSender}.}
	\begin{center}
		\begin{tabularx}{\linewidth}{@{}>{\raggedright\arraybackslash}p{3.6cm} X@{}}
			\toprule
			\textbf{Item} & \textbf{Role in the artifact} \\
			\midrule
			\texttt{sendInterval}, \texttt{payloadBytes} &
			  Periodic offered load; overridden by \texttt{\_netload\_*} overlays. \\
			\addlinespace
			\texttt{dscp} (default~0) &
			  Best Effort marking carried as \texttt{DscpReq}. \\
			\addlinespace
			\texttt{localPort}/\texttt{destPort} (9001) &
			  UDP endpoints for BE TX and for all BE/VO RX KPIs. \\
			\addlinespace
			\texttt{voDedupWindow} &
			  Enables run-long Voice deduplication keyed by (source,~sequence); the time-valued name is retained for legacy configs, but entries are not mid-run expired. \\
			\addlinespace
			Statistics &
			  \texttt{beTxPackets}, \texttt{beRxPackets}, \texttt{voRxPackets}, \texttt{beEndToEndDelay}, \texttt{voEndToEndDelay} (vector + count/mean/min/max). \\
			\bottomrule
		\end{tabularx}
	\end{center}
	\fonte{Prepared by the author.}
\end{table}

On transmission, the module stamps \texttt{CreationTimeTag} and emits a Best Effort counter.
On reception it ignores datagrams whose Layer-3 source equals the local \texttt{wlan0} address, maps remaining packets by DSCP into BE or VO delay/RX streams, and---when deduplication is enabled---accepts only the first Voice packet of each (source,~sequence) pair for the remainder of the run.
This design keeps a single receiver module as the KPI sink for both traffic classes.

\subsection{\texttt{CrashBurstApp}: crash-triggered Voice}
\label{sec:impl:crash}

\texttt{veins\_qos.traffic.CrashBurstApp} implements the crash-alert overlay~\cite{RFC3246_EF_PHB,aguiar2026_veins_qos}.
Every vehicle instantiates the module, but only the node whose OMNeT++ index equals \texttt{targetNodeIndex} (default~0) schedules crash behaviour; the sentinel value~$-1$ would activate all nodes and is unused in the active matrix.

At absolute simulation time \texttt{crashAt} (30~s in the highway packages), the target node:

\begin{enumerate}
  \item marks itself visually (red icon) for interactive inspection;
  \item commands SUMO through TraCI to set speed to zero, realizing the stationary crash source of \autoref{sec:sys:event}~\cite{veins_sommer2011,lopez2018sumo};
  \item begins emitting Voice bursts marked with DSCP~46, each logical sequence comprising \texttt{repeatCount} physical packets spaced by \texttt{repeatGap} with optional \texttt{repeatJitter}.
\end{enumerate}

After \texttt{resumeAfter} (30~s), Voice transmission stops and TraCI restores normal speed control (\texttt{setSpeed(-1)}).
The crash source continues to run \texttt{CritPacketSender} throughout, so alerts are superimposed on ongoing Best Effort traffic rather than replacing it (\autoref{sec:sys:traffic}).

\autoref{tab:impl:crash} summarizes the crash-application surface.
Physical Voice transmissions increment \texttt{voTxPackets}; each logical burst sequence increments \texttt{voLogicalTxPackets}.
The evaluation prefers logical TX when present so that repeatCount does not inflate the denominator of reception ratios; the implementation exposes both counters precisely for that reason~\cite{aguiar2026_veins_qos}.

\begin{table}[htb]
	\ABNTEXfontereduzida
	\caption{\label{tab:impl:crash}Principal parameters and statistics of \texttt{CrashBurstApp}.}
	\begin{center}
		\begin{tabularx}{\linewidth}{@{}>{\raggedright\arraybackslash}p{3.6cm} X@{}}
			\toprule
			\textbf{Item} & \textbf{Role in the artifact} \\
			\midrule
			\texttt{targetNodeIndex} &
			  Which OMNeT++ node becomes the crash source (default~0). \\
			\addlinespace
			\texttt{crashAt}, \texttt{resumeAfter} &
			  Absolute crash start and crash-window duration. \\
			\addlinespace
			\texttt{sendInterval}, \texttt{payloadBytes} &
			  Inter-burst period and payload during the window. \\
			\addlinespace
			\texttt{repeatCount}, \texttt{repeatGap}, \texttt{repeatJitter} &
			  Physical repeats of each logical alert. \\
			\addlinespace
			\texttt{localPort}~9002 / \texttt{destPort}~9001 &
			  TX on a dedicated local port; RX KPIs collected by Crit on~9001. \\
			\addlinespace
			Statistics &
			  \texttt{voTxPackets} (physical), \texttt{voLogicalTxPackets} (per sequence). \\
			\bottomrule
		\end{tabularx}
	\end{center}
	\fonte{Prepared by the author.}
\end{table}

There is no multi-level criticality ladder and no separate four-priority application family: the implementation stays on the two-class \gls{be}/\gls{vo} abstraction fixed in the System Model.

%%%%%%%%%%%%%%%%%%%%%%%%%%%%%%%%%%%%%%%%%%%%%%%%%%%%%%%%%%%%%%%%%%%%%%%%%%%%%%%%
\section{QoS Classifier}
\label{sec:impl:classifier}

Correct placement of crash alerts into the Voice access category is not automatic under the common three-MSB DSCP-to-\gls{up} heuristic (\autoref{sec:bg:rfc8325}, \autoref{sec:sys:qos}).
The artifact therefore installs an explicit classifier rather than relying on default DiffServ-to-802.11 mapping~\cite{RFC8325_WiFi_QoS}.

\texttt{veins\_qos.qos.QosClassifier} implements INET's \texttt{IIeee8021dQosClassifier} interface.
It reads DSCP from packet indications or requests (and, if needed, from IPv4/IPv6 headers), maps \texttt{crashDscp} (default~46) to Voice user priority, and assigns \texttt{defaultUp} (default~0, Best Effort) to every other packet~\cite{RFC8325_WiFi_QoS,IEEE80211e_2005,aguiar2026_veins_qos}.
The resulting \texttt{UserPriorityReq} selects the corresponding \gls{edca} access category inside the \gls{hcf}.

The classifier is installed only under EDCA policies:

\begin{quote}
\ttfamily
*.node[*].wlan[*].classifier.typename = "veins\_qos.qos.QosClassifier"\\
*.node[*].wlan[*].classifier.crashDscp = 46\\
*.node[*].wlan[*].classifier.defaultUp = 0
\end{quote}

Under \texttt{plain}, no classifier is configured: a single \gls{dcf} pending queue carries all traffic without DSCP-to-\gls{ac} differentiation, matching the non-QoS baseline of \autoref{sec:sys:qos}.

%%%%%%%%%%%%%%%%%%%%%%%%%%%%%%%%%%%%%%%%%%%%%%%%%%%%%%%%%%%%%%%%%%%%%%%%%%%%%%%%
\section{MAC Layer: Instrumented MAC and Adaptive HCF}
\label{sec:impl:mac}

All active configurations replace the stock IEEE~802.11 \gls{mac} typename with an instrumented subclass, while only the adaptive policies replace the \gls{hcf} submodule~\cite{IEEE80211e_2005,inet_framework,aguiar2026_veins_qos}.

\subsection{\texttt{V2xIeee80211Mac}}
\label{sec:impl:v2xmac}

\texttt{veins\_qos.mac.V2xIeee80211Mac} extends INET's \texttt{Ieee80211Mac} and is selected globally:

\begin{quote}
\ttfamily *.node[*].wlan[*].mac.typename = "veins\_qos.mac.V2xIeee80211Mac"
\end{quote}

The module subscribes to \texttt{packetDroppedSignal} and, at finish time, records scalar drop totals partitioned by access category (BK, BE, VI, VO, and unclassified) and by drop reason (queue overflow, retry-limit reached, congestion, incorrectly received, and related INET reasons).
Those counters are the \gls{mac}-side counterpart to the application delay and reception statistics: they expose where frames were discarded under each policy without requiring Evaluation-side reinterpretation of raw event logs~\cite{aguiar2026_veins_qos}.

\subsection{\texttt{V2xHcf} and \texttt{V2xEdcaFsmController}}
\label{sec:impl:v2xhcf}

The adaptive policies substitute

\begin{quote}
\ttfamily *.node[*].wlan[*].mac.hcf.typename = "veins\_qos.mac.V2xHcf"
\end{quote}

for the stock INET \texttt{Hcf}.
\texttt{V2xHcf} implements the \texttt{IHcf} contract and embeds a full QoS \gls{hcf} submodule tree (EDCA, HCCA placeholder, originator/recipient QoS data services, rate selection, acknowledgement and RTS/CTS policies) plus a study-specific \texttt{FSMController} of type \texttt{V2xEdcaFsmController}~\cite{IEEE80211e_2005,aguiar2026_veins_qos}.
Because only the typename changes, queue depths and contention parameters inherited from \texttt{edca\_only} remain numerically identical across the adaptive profiles.

\autoref{tab:impl:v2xhcf} lists the controller-facing parameters declared on \texttt{V2xHcf} and its FSM submodule.
Their semantics match the three-mode controller of \autoref{sec:sys:controller}; this section records how those knobs are named in software.

\begin{table}[htb]
	\ABNTEXfontereduzida
	\caption{\label{tab:impl:v2xhcf}Adaptive controller parameters exposed by \texttt{V2xHcf} / \texttt{V2xEdcaFsmController}.}
	\begin{center}
		\begin{tabularx}{\linewidth}{@{}>{\ttfamily\raggedright\arraybackslash}p{4.0cm} X@{}}
			\toprule
			\textbf{Parameter} & \textbf{Meaning} \\
			\midrule
			adaptiveBlocking &
			  Enables Voice-driven Best Effort suppression through the FSM. \\
			\addlinespace
			emergencyPreemption &
			  When true, newly arriving Best Effort may be dropped and stale Best Effort grants suppressed while protection is active. \\
			\addlinespace
			blockDuration &
			  Protection-window extension per Voice demand event (local enqueue or overheard Voice frame). \\
			\addlinespace
			maxContinuousBlock &
			  Cap on one continuous blocking/sending interval (also mirrored on the FSM). \\
			\addlinespace
			voQueueThreshold &
			  Local Voice queue depth that can trigger protection from upper traffic. \\
			\addlinespace
			FSMController.sendingGuardTimeout &
			  Short guard that delays return to listening around Voice send transitions. \\
			\bottomrule
		\end{tabularx}
	\end{center}
	\fonte{Prepared by the author.}
\end{table}

\subsection{Controller integration hooks}
\label{sec:impl:hooks}

\autoref{tab:impl:hooks} summarizes how \texttt{V2xHcf} couples application- and peer-observed Voice activity to Best Effort eligibility.
The controller does not rewrite backoff arithmetic or introduce scheduled gates; it interposes on channel-access eligibility for Best Effort while Voice frames remain subject to ordinary \gls{edca} contention~\cite{IEEE80211e_2005}.

\begin{table}[htb]
	\ABNTEXfontereduzida
	\caption{\label{tab:impl:hooks}Principal integration hooks of \texttt{V2xHcf} relative to the stock INET \gls{hcf}.}
	\begin{center}
		\begin{tabularx}{\linewidth}{@{}>{\raggedright\arraybackslash}p{3.4cm} X@{}}
			\toprule
			\textbf{Hook} & \textbf{Behaviour in the adaptive policies} \\
			\midrule
			Voice demand detection &
			  Local Voice enqueue above \texttt{voQueueThreshold}, or an overheard Voice-classified \gls{mac} frame, calls into the FSM (\texttt{onVoDemandDetected}) and may start or extend a protection window of length \texttt{blockDuration}, capped by \texttt{maxContinuousBlock}. \\
			\addlinespace
			Upper-frame gating &
			  While the FSM reports Best Effort blocked (\emph{blocking} or \emph{sending}), Best Effort channel-access requests are deferred. \\
			\addlinespace
			Own Voice on air &
			  Start and end of the node's own Voice transmissions drive \texttt{onVoTransmissionStart}/\texttt{End}, moving the FSM through \emph{sending} and applying \texttt{sendingGuardTimeout} before listening resumes. \\
			\addlinespace
			Emergency preemption &
			  If \texttt{emergencyPreemption} is true, newly arriving Best Effort frames may be dropped with a congestion reason and stale Best Effort grants suppressed while protection is active. \\
			\addlinespace
			Instrumentation &
			  Signals/scalars \texttt{beDroppedWhileBlocked}, \texttt{beGrantSuppressedWhileBlocked}, and \texttt{voProtectionActivation}; FSM vector \texttt{v2xState} with encoding 0~=~listening, 1~=~blocking, 2~=~sending. \\
			\bottomrule
		\end{tabularx}
	\end{center}
	\fonte{Prepared by the author.}
\end{table}

%%%%%%%%%%%%%%%%%%%%%%%%%%%%%%%%%%%%%%%%%%%%%%%%%%%%%%%%%%%%%%%%%%%%%%%%%%%%%%%%
\section{MAC Policy Configuration Wiring}
\label{sec:impl:policies}

\autoref{tab:impl:map} restates the stack mapping.
The subsections below document the exact OMNeT++ configuration composition used in the active highway \texttt{omnetpp.ini} files~\cite{aguiar2026_veins_qos}.
Numeric access-control values are collected in \autoref{tab:impl:parameters}; here the emphasis is on \emph{which} typename and parameter paths bind them.

\begin{table}[htb]
	\ABNTEXfontereduzida
	\caption{\label{tab:impl:map}Artifact mapping per stack layer.}
	\begin{center}
		\begin{tabularx}{\linewidth}{@{}l >{\raggedright\arraybackslash}X >{\raggedright\arraybackslash}X@{}}
			\toprule
			\textbf{Layer} & \textbf{Reused baseline} & \textbf{Introduced modules} \\
			\midrule
			App. &
			  Multicast UDP hooks in \texttt{VeinsInetApplicationBase} &
			  \texttt{CritPacketSender}, \texttt{CrashBurstApp}; KPI tags and deduplication \\
			\addlinespace
			Net/\gls{qos} &
			  IPv4 forwarding and transport &
			  \texttt{QosClassifier} (\gls{dscp}~46~$\rightarrow$ Voice \gls{up}) \\
			\addlinespace
			\gls{mac} &
			  IEEE~802.11p \gls{dcf}/\gls{edca} queues and channel access &
			  \texttt{V2xIeee80211Mac}; \texttt{V2xHcf}+\texttt{V2xEdcaFsmController} on adaptive policies \\
			\addlinespace
			\gls{phy} &
			  Dimensional radio, propagation, mobility, TraCI coupling &
			  None (shared across all policies) \\
			\bottomrule
		\end{tabularx}
	\end{center}
	\fonte{Prepared by the author.}
\end{table}

\subsection{\texttt{plain}}
\label{sec:impl:plain}

\texttt{[Config plain]} keeps a non-QoS station.
It sets the single DCF pending-queue capacity to 128~packets and does not enable \texttt{qosStation}, does not install \texttt{QosClassifier}, and does not override \texttt{hcf.typename}.
DSCP markings still travel with IP packets for KPI classification at receivers, but they do not drive EDCA access categories.

\subsection{\texttt{edca\_only}}
\label{sec:impl:edca}

\texttt{[Config edca\_only]} extends the plain baseline with standard EDCA~\cite{IEEE80211e_2005}:

\begin{itemize}
  \item \texttt{mac.qosStation = true};
  \item \texttt{mac.hcf.typename = "Hcf"} (stock INET);
  \item \texttt{QosClassifier} as above;
  \item pending-queue capacities BK/BE/VI/VO = 128/128/64/32 (\texttt{edcaf[0..3]});
  \item MSDU aggregation policy cleared (\texttt{msduAggregationPolicy.typename = ""});
  \item \gls{aifsn} BK/BE/VI/VO = 7/3/2/2;
  \item \texttt{cwMin}/\texttt{cwMax} BK/BE = 15--1023, VI = 7--15, VO = 3--7.
\end{itemize}

INET indexes \texttt{edcaf[0]=BK}, \texttt{[1]=BE}, \texttt{[2]=VI}, \texttt{[3]=VO}; the active applications exercise only BE and VO, so VI/BK queues exist but remain empty under the two-class traffic model.

\subsection{Adaptive profiles}
\label{sec:impl:adaptive}

Each adaptive configuration extends \texttt{edca\_only}, replaces \texttt{hcf.typename} with \texttt{veins\_qos.mac.V2xHcf}, and sets \texttt{adaptiveBlocking = true}.
\autoref{tab:impl:policyknobs} records the profile-specific knobs exactly as bound in \texttt{omnetpp.ini}.

\begin{table}[htb]
	\ABNTEXfontereduzida
	\caption{\label{tab:impl:policyknobs}Adaptive policy knobs as bound in the highway \texttt{omnetpp.ini}.}
	\begin{center}
		\begin{tabularx}{\linewidth}{@{}l >{\raggedright\arraybackslash}X@{}}
			\toprule
			\textbf{Config} & \textbf{Controller binding} \\
			\midrule
			\texttt{edca\_v2x\_vo\_stable} &
			  \texttt{blockDuration=15ms}, \texttt{maxContinuousBlock=80ms}, \texttt{sendingGuardTimeout=5ms}, \texttt{voQueueThreshold=2}; \texttt{emergencyPreemption} left false. \\
			\addlinespace
			\texttt{edca\_v2x\_vo\_guarded} &
			  \texttt{blockDuration=4ms}, \texttt{maxContinuousBlock=20ms}, \texttt{sendingGuardTimeout=4ms}, \texttt{voQueueThreshold=3}; preemption false. \\
			\addlinespace
			\texttt{edca\_v2x\_vo\_emergency} &
			  \texttt{emergencyPreemption=true}, \texttt{blockDuration=10ms}, \texttt{maxContinuousBlock=60ms}, \texttt{sendingGuardTimeout=5ms}, \texttt{voQueueThreshold=1}. \\
			\bottomrule
		\end{tabularx}
	\end{center}
	\fonte{Prepared by the author.}
\end{table}

Legacy configuration names from earlier prototypes (for example Friendly/Protect aliases documented in the artifact's deprecation notes) are not part of the active matrix and must not be confused with the five policies of \autoref{sec:sys:qos}.

%%%%%%%%%%%%%%%%%%%%%%%%%%%%%%%%%%%%%%%%%%%%%%%%%%%%%%%%%%%%%%%%%%%%%%%%%%%%%%%%
\section{Scenario Packages, Simulation Parameters, and Configuration Matrix}
\label{sec:impl:scenarios}

\subsection{Consolidated simulation parameters}
\label{sec:impl:parameters}

\autoref{tab:impl:parameters} collects the concrete timing, topology, radio, application, classifier, and \gls{mac} bindings used by the active highway artifact.
It is the numeric counterpart of the System Model in \autoref{chap:sysmodel}: the conceptual roles of the five policies and the crash timeline are defined there; the values that instantiate them in this study are recorded here~\cite{aguiar2026_veins_qos}.
Best Effort inter-arrivals use OMNeT++ \texttt{exponential($\cdot$)} draws with the stated mean; Voice burst periods are fixed, with optional per-repeat jitter.
Adaptive policies are abbreviated \texttt{stable}, \texttt{guarded}, and \texttt{emergency} in the first column; the full OMNeT++ configuration names are \texttt{edca\_v2x\_vo\_*}.

\begin{table}[H]
	\ABNTEXfontereduzida
	\caption{\label{tab:impl:parameters}Simulation parameters used in the crash-aware Veins/\gls{inet} artifact.}
	\begin{center}
	\setlength{\tabcolsep}{3.5pt}
	\begin{tabularx}{\linewidth}{@{}
		>{\raggedright\arraybackslash}p{2.6cm}
		>{\raggedright\arraybackslash}p{2.9cm}
		>{\raggedright\arraybackslash}X
		@{}}
		\toprule
		\multicolumn{3}{@{}l@{}}{\textbf{Common highway setting}} \\
		\midrule
		\textbf{Parameter} & \textbf{Value} & \textbf{Notes} \\
		\midrule
		Sim.\ time &
		  \SI{70}{\second} &
		  \texttt{sim-time-limit}; identical for every configuration. \\
		\addlinespace
		Corridor &
		  5~km, 3+3 lanes &
		  Straight bidirectional SUMO corridor, three lanes per direction; mixed car/truck types (EIDM, LC2013); nominal lane speed \SI{120}{\km\per\hour}. \\
		\addlinespace
		Density &
		  10~/~100~veh. &
		  Light and heavy packages (\texttt{veins\_inet\_highway\_light}/\texttt{\_heavy}); same \texttt{omnetpp.ini}, different \texttt{highway.rou.xml}. \\
		\addlinespace
		Crash &
		  \SI{30}{\second}--\SI{60}{\second} &
		  Node index~0 (\texttt{node0\_highway}): TraCI stop, DSCP~46 alerts; BE background continues. \\
		\addlinespace
		Radio &
		  802.11p / ch.~3 &
		  5.9~GHz, 10~MHz, 20~mW; \texttt{Ieee80211DimensionalRadio}; mode~\texttt{p}; data frames at 12~Mbit/s (fastest mandatory mode, \gls{inet} default rate selection). \\
		\addlinespace
		Propagation &
		  FS~+~binary &
		  Free-space path loss; no fading/shadowing; \texttt{IdealObstacleLoss} + \texttt{highway.obstacles.xml}; SNIR~\SI{4}{\decibel}; sens.\ $-$85~dBm; noise $-$110~dBm (\gls{inet}~4.5.4 defaults, not overridden). \\
		\addlinespace
		Multicast &
		  \texttt{224.0.0.1} &
		  On \texttt{wlan0}. Crit ports 9001/9001; Crash local~9002, dest~9001 (VO RX counted in Crit). \\
		\addlinespace
		TraCI &
		  port~9999 &
		  \texttt{localhost}; \texttt{VeinsInetManager} + \texttt{veins\_launchd}; update~\SI{1}{\second}; SUMO step~\SI{0.1}{\second}. \\
		\addlinespace
		Host/\gls{mac} &
		  \texttt{VeinsInetCar} &
		  Always \texttt{V2xIeee80211Mac}; adaptive policies replace \texttt{Hcf} with \texttt{V2xHcf}. \\
		\addlinespace
		Classifier &
		  DSCP~46~$\rightarrow$~VO &
		  \texttt{QosClassifier} on EDCA policies only (\texttt{crashDscp=46}, \texttt{defaultUp=0}); unused under \texttt{plain}. \\
		\addlinespace
		Recording &
		  vec.\ +\ sca. &
		  \texttt{**.vector-recording}/\texttt{scalar-recording = true}. \\
		\addlinespace
		Matrix &
		  15~/~package &
		  Five policies $\times$ three loads; $\times$ two densities $=$ 30 named experiments before seeds. \\
		\midrule
		\multicolumn{3}{@{}l@{}}{\textbf{Application load profiles (\texttt{\_netload\_*})}} \\
		\midrule
		\textbf{Profile} & \textbf{BE (\texttt{app[0]})} & \textbf{VO crash (\texttt{app[1]})} \\
		\midrule
		\texttt{low} &
		  exp.~\SI{500}{\milli\second}, 200~B &
		  \SI{120}{\milli\second}, 150~B; repeat~3; gap~\SI{20}{\milli\second}; jitter~\SI{5}{\milli\second}. \\
		\addlinespace
		\texttt{medium} &
		  exp.~\SI{250}{\milli\second}, 320~B &
		  \SI{75}{\milli\second}, 180~B; repeat~4; gap~\SI{15}{\milli\second}; jitter~\SI{5}{\milli\second}. \\
		\addlinespace
		\texttt{high} &
		  exp.~\SI{125}{\milli\second}, 420~B &
		  \SI{20}{\milli\second}, 260~B; repeat~8; gap~\SI{5}{\milli\second}; jitter~\SI{2}{\milli\second}. \\
		\midrule
		\multicolumn{3}{@{}l@{}}{\textbf{\gls{mac} policies (runnable name \texttt{<policy>\_netload\_<profile>})}} \\
		\midrule
		\textbf{Policy} & \textbf{Role} & \textbf{Access-control parameters} \\
		\midrule
		\texttt{plain} &
		  Non-\gls{qos} \gls{dcf} &
		  Single pending queue~128; no classifier; no \gls{edca}; no adaptive blocking. \\
		\addlinespace
		\texttt{edca\_only} &
		  Stock EDCA &
		  \texttt{Hcf} + \texttt{QosClassifier}. DSCP~46~$\rightarrow$~\texttt{AC\_VO}, else UP~0. Queues BK/BE/VI/VO: 128/128/64/32. \gls{aifsn} 7/3/2/2. \texttt{cwMin} 15/15/7/3. \texttt{cwMax} 1023/1023/15/7. \\
		\addlinespace
		\texttt{stable} &
		  Sustained VO &
		  Config \texttt{edca\_v2x\_vo\_stable}. Inherits \texttt{edca\_only}; \texttt{V2xHcf}. Block~\SI{15}{\milli\second}; max cont.~\SI{80}{\milli\second}; guard~\SI{5}{\milli\second}; VO thr.~2. \\
		\addlinespace
		\texttt{guarded} &
		  Bounded VO &
		  Config \texttt{edca\_v2x\_vo\_guarded}. Same inheritance. Block~\SI{4}{\milli\second}; max cont.~\SI{20}{\milli\second}; guard~\SI{4}{\milli\second}; VO thr.~3. \\
		\addlinespace
		\texttt{emergency} &
		  Emergency VO &
		  Config \texttt{edca\_v2x\_vo\_emergency}; \texttt{emergencyPreemption=true}. Block~\SI{10}{\milli\second}; max cont.~\SI{60}{\milli\second}; guard~\SI{5}{\milli\second}; VO thr.~1; may drop BE. \\
		\bottomrule
	\end{tabularx}
	\end{center}
	\fonte{Prepared by the author.}
\end{table}

\subsection{Density packages}
\label{sec:impl:density}

The dissertation experiments use two parallel packages under
\texttt{veins\_qos/simulations/}~\cite{aguiar2026_veins_qos}:

\begin{itemize}
  \item \texttt{veins\_inet\_highway\_light} --- light-density corridor with 10~vehicles;
  \item \texttt{veins\_inet\_highway\_heavy} --- heavy-density corridor with 100~vehicles.
\end{itemize}

Both packages share the same corridor geometry (\texttt{highway.net.xml}), obstacle file, launchd description, and---byte-for-byte in the active tree---the same \texttt{omnetpp.ini}.
Fleet size is selected by the SUMO route file (\texttt{highway.rou.xml}): a fixed vehicle \texttt{node0\_highway} departs at time~0 and becomes OMNeT++ node index~0 (the default crash target), while additional vehicles are injected from SUMO flows.
Separating density (package) from offered load (ini overlay) is deliberate: the two stress dimensions can be varied independently~\cite{aguiar2026_veins_qos}.

\subsection{Road network and mobility}
\label{sec:impl:sumo}

The SUMO network is a straight bidirectional highway corridor of 5~km with three lanes per direction and mixed car/truck vehicle types under EIDM car-following and LC2013 lane-changing models~\cite{lopez2018sumo}.
Simulation step size is 0.1~s; teleportation on gridlock is disabled (\texttt{time-to-teleport=-1}) so that stopped crash nodes remain in place.
Veins updates network-node mobility from TraCI; the manager polls on a one-second update interval in the highway ini, while SUMO itself advances at its finer step~\cite{veins_sommer2011}.

\subsection{Radio and propagation baseline}
\label{sec:impl:radio}

Every configuration shares the IEEE~802.11p radio profile summarized in \autoref{tab:impl:parameters}~\cite{IEEE80211,inet_framework}:

\begin{itemize}
  \item operation mode \texttt{p}, band 5.9~GHz, channel~3, bandwidth 10~MHz, transmit power 20~mW;
  \item radio typename \texttt{Ieee80211DimensionalRadio};
  \item antenna mobility attached to the vehicle mobility module with a fixed longitudinal/height offset;
  \item obstacle loss typename \texttt{IdealObstacleLoss}, configured from
        \texttt{highway.obstacles.xml}~\cite{inet_ideal_obstacle_loss};
  \item free-space path loss and a static SNIR threshold, without small-scale fading or shadowing (\autoref{sec:sys:scope}).
\end{itemize}

Receiver sensitivity ($-$85~dBm), energy detection ($-$85~dBm), the \gls{snir} threshold (\SI{4}{\decibel}), background noise ($-$110~dBm), free-space path loss, and constant-speed propagation are the \gls{inet}~4.5.4 defaults for this radio and medium; they are not overridden in the configuration.
No rate-control module is configured, so \gls{inet}'s default rate selection transmits data frames---including group-addressed ones---at the fastest mandatory mode of the \SI{10}{\mega\hertz} \gls{ofdm} mode set (12~Mbit/s).
This differs from production deployments, where group-addressed frames are commonly sent at a basic rate (\autoref{sec:bg:multicast}); the choice is uniform across all five policies.

Two fidelity caveats apply.
First, the artifact realizes 802.11p-like operation through \gls{inet}'s \texttt{AdhocHost} (no association or authentication, ad hoc management); it approximates \gls{ocb} operation rather than literally implementing \texttt{dot11OCBActivated} behaviour.
Second, the \gls{edca} parameter set of \autoref{tab:impl:parameters} follows the IEEE~802.11 BSS QoS defaults, whereas IEEE~802.11 specifies a distinct default \gls{edca} parameter set for \gls{ocb} operation with a less aggressive Best Effort category (\gls{aifsn}~6 instead of~3)~\cite{IEEE80211,etsi_en_302663}.
The configured set therefore makes Best Effort compete harder than an \gls{ocb}-default deployment would, which is conservative for the protection problem studied here.

Because this baseline is identical across policies, cross-policy differences isolate queueing, contention, and local Best Effort suppression.

\subsection{Timing and default application binding}
\label{sec:impl:timing}

Under \texttt{[General]}, both packages set \texttt{sim-time-limit = 70s}, crash start \texttt{crashAt = 30s}, and crash duration \texttt{resumeAfter = 30s}, matching \autoref{sec:sys:event} and \autoref{tab:impl:parameters}.
Default application parameters in \texttt{[General]} match the medium load profile; the \texttt{\_netload\_*} overlays override them for the full matrix.

\subsection{Load overlays and runnable matrix}
\label{sec:impl:matrix}

Three helper configurations vary only application-side load, using the means and payloads in \autoref{tab:impl:parameters}.
Runnable experiments compose a policy configuration with a load overlay via OMNeT++ \texttt{extends}, for example

\begin{quote}
\ttfamily
[Config edca\_v2x\_vo\_emergency\_netload\_high]\\
extends = edca\_v2x\_vo\_emergency, \_netload\_high
\end{quote}

yielding the naming pattern \texttt{<policy>\_netload\_<low|medium|high>}: five policies $\times$ three loads $=$ fifteen named configurations per density package~\cite{aguiar2026_veins_qos}.
Across both light and heavy packages the unique named experiments therefore total thirty before seed repetitions (seed counts belong to Evaluation).

Batch execution helpers in each package (\texttt{run}, \texttt{run\_matrix.sh}) select profiles such as high-load-only cores or the full fifteen-config matrix; they presuppose a running \texttt{veins\_launchd} on port~9999.
Building the study library follows the standard OMNeT++ pattern (\texttt{make makefiles \&\& make} under \texttt{veins\_qos/}) with NED paths that include INET, Veins, and \texttt{veins\_inet} sources~\cite{aguiar2026_veins_qos}.

%%%%%%%%%%%%%%%%%%%%%%%%%%%%%%%%%%%%%%%%%%%%%%%%%%%%%%%%%%%%%%%%%%%%%%%%%%%%%%%%
\section{Instrumentation and Analysis Surface}
\label{sec:impl:instrumentation}

\subsection{Recorded signals}
\label{sec:impl:signals}

Highway configurations enable comprehensive recording:

\begin{quote}
\ttfamily **.vector-recording = true\\
**.scalar-recording = true
\end{quote}

\autoref{tab:impl:signals} collects the study-relevant statistics emitted by the modules described above.
These signals are the measurement surface consumed by Evaluation; their statistical treatment (seeds, pooling, confidence reporting) is deferred to that chapter.

\begin{table}[htb]
	\ABNTEXfontereduzida
	\caption{\label{tab:impl:signals}Study-relevant instrumentation exposed by the artifact.}
	\begin{center}
		\begin{tabularx}{\linewidth}{@{}l >{\raggedright\arraybackslash}X@{}}
			\toprule
			\textbf{Source} & \textbf{Recorded quantities} \\
			\midrule
			\texttt{CritPacketSender} &
			  BE TX/RX counts; VO RX counts; BE and VO end-to-end delay vectors with count/mean/min/max scalars. \\
			\addlinespace
			\texttt{CrashBurstApp} &
			  Physical VO TX count; logical VO sequence count. \\
			\addlinespace
			\texttt{V2xIeee80211Mac} &
			  Per-\gls{ac} drop totals; per-reason drop totals (overflow, retry limit, congestion, \ldots). \\
			\addlinespace
			\texttt{V2xHcf} &
			  Best Effort drops while blocked; Best Effort grant suppressions; Voice protection activations. \\
			\addlinespace
			\texttt{V2xEdcaFsmController} &
			  \texttt{v2xState} vector (listening/blocking/sending). \\
			\bottomrule
		\end{tabularx}
	\end{center}
	\fonte{Prepared by the author.}
\end{table}

End-to-end delay is computed at receivers as simulation time minus \texttt{CreationTimeTag}.
For crash frames the creation stamp is the actual send instant of each physical repeat, so delay reflects channel and queueing latency rather than logical inter-burst spacing~\cite{aguiar2026_veins_qos}.

\subsection{Offline analysis tooling}
\label{sec:impl:kpi}

The companion \texttt{kpi\_dashboard/} project parses OMNeT++ scalar and vector result files, caches intermediate tables, and exports publication figures~\cite{aguiar2026_veins_qos}.
From an implementation perspective its contract is:

\begin{itemize}
  \item ingest \texttt{.sca}/\texttt{.vec} outputs under each scenario's \texttt{results/} directory;
  \item aggregate application delays and counts from \texttt{Scenario.node[*].app[0]} (and crash TX from \texttt{app[1]});
  \item expose controller and drop counters when present under adaptive policies;
  \item write tabular and figure artifacts consumed by Evaluation.
\end{itemize}

The dashboard does not alter simulation behaviour; it is a reproducibility aid that turns the instrumentation of \autoref{tab:impl:signals} into analysis-ready products.
Interpretation of those products---protection-versus-cost trade-offs, confidence intervals, and discussion---belongs to the Evaluation chapter.

%%%%%%%%%%%%%%%%%%%%%%%%%%%%%%%%%%%%%%%%%%%%%%%%%%%%%%%%%%%%%%%%%%%%%%%%%%%%%%%%
\section{Implementation Scope and Deliberate Omissions}
\label{sec:impl:scope}

The code base stays aligned with the minimal System Model of \autoref{sec:sys:scope}:

\begin{itemize}
  \item a single crash source and a two-class \gls{be}/\gls{vo} distinction only;
  \item no infrastructure, roadside unit, or Time-Sensitive Networking-style gating;
  \item no \gls{dcc} rate control---application rates follow the fixed offered-load overlays;
  \item no multi-crash orchestration and no standardized CAM/DENM stack;
  \item no small-scale fading, shadowing, or graded obstacle attenuation beyond binary \texttt{IdealObstacleLoss}~\cite{inet_ideal_obstacle_loss}.
\end{itemize}

Older packages under \texttt{simulations/veins\_inet\_square/},
\texttt{veins\_inet\_highway/}, and \texttt{veins\_inet\_light/}
are retained for development history and must not be mixed with the active light/heavy highway result trees when analysing the dissertation campaigns~\cite{aguiar2026_veins_qos}.

Together, the toolchain, host compounds, packet path, applications, classifier, instrumented \gls{mac}, adaptive \gls{hcf}, policy wiring, scenario packages, and instrumentation constitute the concrete realization of the crash-aware System Model.
\autoref{chap:evaluation} asks whether crash-triggered Voice traffic obtains better wireless service than ordinary Best Effort under contention, and what cost that protection imposes on ordinary traffic.
